\documentclass[11pt,a4paper]{article}
\usepackage[T1]{fontenc}
\usepackage[utf8]{inputenc}
\usepackage{lmodern,microtype}
\usepackage[margin=25mm,headheight=14pt]{geometry}
\usepackage{amsmath,amssymb,amsthm,mathtools}
\usepackage{booktabs,array,longtable}
\usepackage{xcolor,listings,fancyhdr}
\usepackage[hidelinks]{hyperref}
\hypersetup{pdftitle={RadicalRoots: Constructive Conversion of Algebraic Numbers to Radicals},pdfsubject={Constructive Galois theory, exact radicals, and a Python/Wolfram implementation}}
\usepackage{enumitem}
\definecolor{ink}{HTML}{203747}
\definecolor{pale}{HTML}{F3F5F6}
\lstset{basicstyle=\ttfamily\small,backgroundcolor=\color{pale},
 breaklines=true,columns=fullflexible,keepspaces=true,showstringspaces=false,
 frame=single,rulecolor=\color{black!15},aboveskip=10pt,belowskip=10pt}
\setlength{\parskip}{5pt}
\setlength{\parindent}{0pt}
\setlist{nosep,leftmargin=*}
\pagestyle{fancy}\fancyhf{}
\fancyhead[L]{\small\textsc{RadicalRoots}}
\fancyhead[R]{\small Constructive radical conversion}
\fancyfoot[C]{\small\thepage}
\newtheorem{theorem}{Theorem}[section]
\newtheorem{proposition}[theorem]{Proposition}
\newtheorem{lemma}[theorem]{Lemma}
\newtheorem{corollary}[theorem]{Corollary}
\theoremstyle{definition}\newtheorem{definition}[theorem]{Definition}
\theoremstyle{remark}\newtheorem{remark}[theorem]{Remark}
\newcommand{\Q}{\mathbb Q}\newcommand{\C}{\mathbb C}
\newcommand{\Gal}{\operatorname{Gal}}\newcommand{\Aut}{\operatorname{Aut}}
\newcommand{\Tr}{\operatorname{Tr}}\newcommand{\Res}{\operatorname{Res}}
\newcommand{\Span}{\operatorname{span}}\newcommand{\ord}{\operatorname{ord}}
\newcommand{\code}[1]{\texttt{#1}}
\newcommand{\inputroot}{\alpha}
\newcommand{\TestCount}{46}
\newcommand{\TestPassed}{46}
\newcommand{\TestPython}{3.13.5}
\newcommand{\TestSymPy}{1.14.0}

\begin{document}
\begin{titlepage}
\vspace*{10mm}
{\color{ink}\Huge\bfseries RadicalRoots\par}
\vspace{4mm}
{\LARGE Constructive conversion of algebraic numbers to radicals\par}
\vspace{8mm}
{\large A Python reference engine, a Wolfram Language front end,\par
and the mathematics behind their guarantees\par}
\vspace{12mm}
Prepared for Vladimir Reshetnikov\\
8 September 2026 \quad\textbar\quad Version 1.0
\vspace{12mm}
\begin{abstract}
An algebraic number has a radical expression precisely when the Galois group
of its irreducible minimal polynomial is solvable. Turning this criterion into
usable code requires considerably more than naming a Galois group: one must
construct radical generators, represent every radicand in an earlier field,
and select the correct complex branches. This report develops an explicit
algorithm using splitting fields, prime cyclic quotients, rational nullspaces,
and exact tower coordinates. A relative-automorphism construction avoids a
second, potentially much larger, factorization after adjoining roots of unity.
Finite trace-projector candidates substantially improve expression size.

The accompanying Python implementation contains this general construction,
structural shortcuts, a non-executable JSON radical format, and a separate
arithmetic-tower checker. The Wolfram Language package calls the engine and
checks the final answer with \code{RootReduce}. The two motivating examples
receive especially short solutions: a reciprocal sextic reduces to a cubic,
and a quintic is a Dickson polynomial. The mathematical algorithm is complete
with exact algebraic primitives and without resource limits. The software is
a reference implementation, not a claim of unrestricted practical performance
or a formally verified theorem prover.
\end{abstract}
\vfill
\textbf{Validation boundary.} The Python suite was executed with SymPy
\TestSymPy{}: \TestPassed{} of \TestCount{} tests passed. The Wolfram front end
and its integration tests are supplied but were not executed: the available
Wolfram service returned an HTTP 404 error, and no local Wolfram kernel was
available. Exact details and a source fingerprint accompany the archive.
\end{titlepage}
\begingroup
\small
\setlength{\parskip}{0pt}
\tableofcontents
\endgroup
\newpage

\section{What problem is being solved?}
The motivating question \cite{question} asks why \code{ToRadicals} can leave
solvable \code{Root} objects unchanged and whether a systematic workaround
exists. Its documentation guarantees conversion for roots of polynomials of
degree at most four, but explicitly notes that some possible radical
representations are not found \cite{toradicals}. Failure of a particular
conversion routine is therefore not a proof of impossibility.

The present input is a \emph{specified exact algebraic number}. The Python
interface receives a univariate polynomial in $\Q[x]$ and a selected root; the
Wolfram interface first computes the rational minimal polynomial of the
supplied exact value. Parameters, approximate coefficients, and roots of
nonpolynomial analytic equations are outside the contract.

\begin{definition}
A radical expression over $\Q$ uses rational constants, arithmetic operations,
and finitely many rational powers. We allow $i$, equivalently $\sqrt{-1}$.
Every noninteger power has its principal complex meaning unless an explicit
root-of-unity multiplier specifies another branch.
\end{definition}
Neither \code{Root}/\code{CRootOf}, an unevaluated inverse function, nor a
trigonometric constant is an acceptable final radical. In particular,
$2\cos(2\pi/11)$ is not the requested output syntax, whereas
$(-1)^{2/11}+(-1)^{-2/11}$ is. The latter expression uses ordinary complex
radicals; it makes no assertion about expressibility using only real radicals.

\subsection{The selected factor is essential}
If $P$ is reducible, solvability must be tested on the minimal polynomial
$f$ of the selected root, not on $P$ as a whole. For example,
\[
 P(x)=(x^5-x-1)(x^2-2)
\]
has a nonsolvable factor, but the selected number $\sqrt2$ certainly has a
radical representation. The implementation removes repeated factors for
indexing, selects a root, and only then obtains its irreducible minimal
polynomial. All conjugates of one irreducible polynomial have the same
radical-solvability status, although their branch multipliers differ.

\subsection{Four logically different outcomes}
\begin{center}
\begin{tabular}{p{.31\linewidth}p{.61\linewidth}}
\toprule
Result & Meaning\\\midrule
\code{Success} & A radical expression has been constructed with exact
algebraic identities and root selection.\\
\code{NotSolvableByRadicals} & The selected irreducible polynomial has a
computed nonsolvable Galois group.\\
\code{NotFound} & The finite shortcuts were explicitly requested and did
not produce an answer. No impossibility claim is made.\\
\code{ResourceLimit} or \code{BackendFailure} & A budget was exhausted or
an algebraic backend operation failed. Solvability remains unresolved.\\
\bottomrule
\end{tabular}
\end{center}
Malformed mathematical input raises a Python \code{ValueError}; the
command-line worker reports \code{InvalidInput}. The Wolfram interface uses
\code{Failure} objects rather than leaving an expression that might be
mistaken for a successful radical conversion.

\section{What the supplied notebook contributes}
The uploaded archive contains \code{FindExtension.nb}. It was inspected as
inert notebook data, not evaluated. A small box-expression reader recovered
220 input cells. The inspection found several complementary exploratory
ideas, not one complete decision-and-construction procedure.

\code{findExtension} and \code{findExtension3} search for lower-degree factor
coefficients by taking polynomial remainders and eliminating unknowns.
\code{project} examines real and imaginary parts of conjugates for degree
reduction. \code{factor} and \code{solve} try to identify the relevant factor
or conjugate. The random \code{fuzz} search forms sums of conjugates and uses
high-precision reconstruction to suggest lower-degree quantities. Finally,
\code{simplify} tries to improve a proposed expression with progressively
more expensive simplification.

These are useful discovery methods. However, a lower-degree intermediate
field need not itself provide a radical generator; a small residual numerical
error does not identify a conjugate exactly; and the notebook's
\code{While[True]} searches do not supply a stopping criterion. Some helpers
also rely on private simplifier names. In particular, the inspected text
contains both \code{Simplify`SimplifyCount} and a spelling
\code{Simplify`SimpifyCount}; the latter deserves attention independently of
any mathematical algorithm.

The new implementation retains the mathematical ideas of intermediate-field
traces and low-degree factorization, but separates discovery from certification.
Its completeness argument does not depend on random sampling, successful
integer-relation reconstruction, a simplification complexity function, or an
undocumented Wolfram function. It does not attempt to reproduce every
notebook experiment or establish that every number appearing there is
solvable. A provenance file records the original notebook's hash and the
inspection method.

\section{The Galois criterion and its constructive content}
Let $f\in\Q[x]$ be the irreducible minimal polynomial of $\alpha$, let
$E\subset\C$ be its splitting field in the chosen complex embedding, and put
$G=\Gal(E/\Q)$.

\begin{theorem}[Galois criterion]
The number $\alpha$ belongs to a radical extension of $\Q$ if and only if
$G$ is solvable.
\end{theorem}
The classical theorem, including cyclic extensions after adjoining roots of
unity, is treated in \cite[Proposition 5.27 and Theorem 5.34]{milne}. For a
single irreducible polynomial, conjugation carries a radical expression for
one root to radical expressions for all its roots. Conversely, a radical
tower can be enlarged by roots of unity and conjugates to a normal tower
with abelian steps; its relevant Galois quotients are solvable. The remainder
of this report makes the sufficient direction computationally explicit.

A group name alone does not construct the answer. Likewise, a tower whose
relative minimal polynomials have degrees $2,3,5,\ldots$ is not automatically
a radical tower: a general degree-five equation is not a fifth-root operation.
The needed extra structure is a \emph{cyclic} extension over a field already
containing the appropriate roots of unity.

\subsection{Choice of roots of unity}
Write
\[
 m=\prod_{p\mid |G|}p,\qquad \zeta=(-1)^{2/m},\qquad
 K=\Q(\zeta),\qquad L=EK.
\]
Only distinct prime factors occur in $m$. This is sufficient because we will
refine every group quotient to prime order, and $K$ contains
$\zeta_p=\zeta^{m/p}$ for every required prime $p$.

Both $E$ and $K$ are normal over $\Q$, so $L$ is normal as well. Restriction
injects $H=\Gal(L/K)$ into $G$; consequently $H$ is solvable whenever $G$ is.
There is no need to assume $E\cap K=\Q$. In fact, failing to account for that
intersection is a common source of wrong degree counts. The implementation
checks
\[
 |H|=[L:K]=\frac{[L:\Q]}{\varphi(m)}.
\]
The rational case $|G|=1$ is handled before this construction.

\section{Concrete fields and their automorphisms}
\subsection{Constructing the splitting field}
Start with $E_0=\Q$. Factor $f$ over $E_j$. If all factors are linear, stop.
Otherwise adjoin a root of a nonlinear factor, preferably a root supplied by
an exact quadratic, cubic, or quartic formula. When such a shortcut is not
available, choose an original root of $f$ not yet contained in $E_j$.

Every adjunction is a root of $f$, lies in the fixed splitting field, and
strictly increases the current field. Once all original roots belong to the
field, the algorithm stops. This is an effective finite procedure using
exact factorization, field membership, and primitive-element construction.
It is not a claim that these operations are cheap.

\subsection{Rational power coordinates}
Represent a number field as
\[
 F=\Q(\theta)\cong\Q[T]/(q(T)),\qquad d=\deg q.
\]
Its power basis $1,\theta,\ldots,\theta^{d-1}$ makes every element a rational
column vector. Field multiplication is polynomial multiplication followed by
reduction modulo $q$. If $a\in F$, write $M_a$ for the rational matrix of
multiplication by $a$.

If $F/\Q$ is normal, factoring $q$ over $F$ gives all $d$ images of $\theta$.
For each image $b$, substitution $\theta\mapsto b$ defines an automorphism.
Composition and equality are computed by exact polynomial substitution in
$F$. The result is a finite multiplication table, with the identity numbered
zero. No numerical recognition of a permutation is used here.

\subsection{A cheaper relative-automorphism construction}
Recomputing all automorphisms by factoring the minimal polynomial of a
primitive element of $L=E(\zeta)$ can be much more expensive than doing so
for $E$. The following construction avoids that second factorization.

Primitive-element computation supplies rational $c_0,c_1$ such that
\[
 \theta=c_0e+c_1\zeta,\qquad E=\Q(e),\qquad L=\Q(\theta).
\]
For every $\sigma\in G$, the only possible extension fixing $\zeta$ sends
$\theta$ to $b_\sigma=c_0\sigma(e)+c_1\zeta$. In the rational power model
of $L$, test the following identities:
\begin{align*}
 q_L(b_\sigma)&=0,\\
 e(b_\sigma)&=\sigma(e),\\
 \zeta(b_\sigma)&=\zeta.
\end{align*}
Here $e(T)$ and $\zeta(T)$ mean their coordinate polynomials in $\theta$.
All three tests take place inside one exact field.

\begin{proposition}
The retained images are exactly the automorphisms of $L$ fixing $K$.
\end{proposition}
\begin{proof}
A root $b_\sigma$ of the irreducible polynomial $q_L$ defines a $\Q$-embedding
$L\to L$, hence an automorphism. The second and third identities establish
the specified action on the two generators. Conversely, any automorphism
fixing $K$ restricts to an element of $G$, and its action on $\theta$ is forced
to be the stated linear combination. Thus no required automorphism is lost.
\end{proof}
The implementation additionally checks the expected group order
$[L:\Q]/\varphi(m)$ and closure of the resulting table. This is an exact
extension test, not an assumption of linear disjointness.

\section{Prime cyclic quotients without subgroup enumeration}
The derived series $G^{(0)}=G$, $G^{(j+1)}=[G^{(j)},G^{(j)}]$ decides
solvability. A nontrivial stationary term proves nonsolvability. This finite
table algorithm is the general fallback; a small-degree SymPy Galois-group
routine is used only as an optional earlier negative decision. Its documented
degree restriction is therefore not a restriction on the general construction
\cite{sympynf}.

For a solvable subgroup $J$, put $D=[J,J]$. Then $J/D$ is finite abelian.
Starting with $D$, greedily adjoin elements while the generated subgroup
remains proper in $J$. When no further element can be adjoined, the result
$J'$ is a maximal proper subgroup containing $D$. Consequently $J'$ is
normal in $J$, and $J/J'$ is a simple finite abelian group, hence cyclic of
prime order. Repeating gives
\[
 H=H_0\triangleright H_1\triangleright\cdots\triangleright H_s=1,
 \qquad H_{i-1}/H_i\cong C_{p_i}.
\]
Normality is required only in the immediately preceding group. Requiring all
$H_i$ to be normal in the original $G$ would impose an unnecessary and
sometimes impossible restriction.

Set $F_i=L^{H_i}$. The corresponding ascending field tower satisfies
\[
 K=F_0\subset F_1\subset\cdots\subset F_s=L,
 \qquad [F_i:F_{i-1}]=p_i.
\]
Each extension $F_i/F_{i-1}$ is cyclic and contains $\zeta_{p_i}$ in its base.

\section{Constructing a radical generator by a nullspace}
Fix adjacent groups $J=H_{i-1}$ and $N=H_i$, with $[J:N]=p$. Choose
$\tau\in J\setminus N$; its restriction generates $\Gal(F_i/F_{i-1})$.
For an automorphism $g$, let $A_g$ be its rational matrix on the power basis
of $L$.

Consider the simultaneous rational linear system
\begin{equation}\label{eq:eigen}
 (A_h-I)v=0\quad(h\text{ in a generating set of }N),\qquad
 (A_\tau-M_{\zeta_p})v=0.
\end{equation}
The corresponding element $r\in L$ is fixed by $N$ and satisfies
$\tau(r)=\zeta_p r$.

\begin{theorem}[The Kummer step]\label{thm:kummer}
The system \eqref{eq:eigen} has a nonzero solution. Every such solution
satisfies
\[
 r^p\in F_{i-1},\qquad F_i=F_{i-1}(r).
\]
\end{theorem}
\begin{proof}
Choose a normal-basis element $b$ for the cyclic extension $F_i/F_{i-1}$.
The sum
\[
 R=\sum_{j=0}^{p-1}\zeta_p^{-j}\tau^j(b)
\]
is nonzero by linear independence of the normal basis and obeys
$\tau(R)=\zeta_pR$. Since $R\in F_i$, it is also fixed by $N$, proving
existence of a nonzero solution to the rational system.

For any nonzero solution $r$, the element $r^p$ is fixed by $N$ and by
$\tau$. Together these generate $J$, so $r^p\in F_{i-1}$. If $r$ belonged
to $F_{i-1}$, then $r=\tau(r)=\zeta_p r$, contradicting $r\ne0$ and
$\zeta_p\ne1$. Since the extension has prime degree, any element outside
$F_{i-1}$ generates $F_i$.
\end{proof}

The key algorithmic point is that no normal basis has to be found: ordinary
rational nullspace computation constructs the required eigenvector. The
multiplication matrix $M_{\zeta_p}$ is essential. Treating $\zeta_p$ as a
rational scalar in a rational matrix equation would be incorrect.

\subsection{Smaller generators by finite trace projections}
A rational nullspace basis can have unattractive denominators. The code also
tries two deterministic alternatives from powers of the original generator
$e$ of $E$. For $b=e$ or $e^2$, when available, form
\[
 T_N(b)=\sum_{h\in N}h(b),\qquad
 R(b)=\sum_{j=0}^{p-1}\zeta_p^{-j}\tau^j(T_N(b)).
\]
A nonzero $R(b)$ lies in the same eigenspace. The program compares the
rational coefficient sizes of $r^p$ for these candidates and for the
nullspace basis vectors, then chooses a small radicand. Explicit fixed-field
and eigenvector identities are checked again for the chosen candidate.

This is a finite expression-size heuristic, not a random search and not a
completeness assumption. The nullspace always supplies a candidate. In the
cyclic-quintic test, this simple change replaces a very large coordinate
formula by a short fifth-root expression displayed in Section~\ref{sec:c5}.

\section{From field generators to a radical expression}
\subsection{Maintaining an exact tower basis}
Initially use the rational basis
\[
 \mathcal B_0=(1,\zeta,\ldots,\zeta^{\varphi(m)-1})
\]
of $K$. Suppose $\mathcal B_{i-1}$ is already represented both as exact
internal field elements and as expressions in earlier radical variables.
Solve the rational coordinate system
\[
 r_i^{p_i}=\sum_{b\in\mathcal B_{i-1}}c_b b.
\]
Its symbolic counterpart is the next radicand $a_i$. Extend the basis by
\[
 \mathcal B_i=\{b r_i^j:
              b\in\mathcal B_{i-1},\ 0\le j<p_i\}.
\]
Theorem~\ref{thm:kummer} proves linear independence and spanning. Thus the
new basis size is exactly $[F_i:\Q]=[L:\Q]/|H_i|$, an invariant checked
by the program.

If $\alpha$ is already fixed by $H_{i-1}$, stop before adding another
radical. Otherwise continue until $H_s=1$. A final rational coordinate solve
expresses $\alpha$ in the accumulated basis. The implementation never hides
an unresolved algebraic number in a coefficient of the final expression.
Internal field generators may still occur in the audit metadata, which is
not the returned radical expression.

\subsection{Choosing the complex branch}
Let $a=r^p\ne0$ in the concrete embedded field. There is a unique
$k\in\{0,\ldots,p-1\}$ such that
\begin{equation}\label{eq:phase}
 r=\zeta_p^k\,a^{1/p},
\end{equation}
where $a^{1/p}$ is the principal complex power. After identifying $k$, emit
that multiplier explicitly. In particular, a principal cube root of a
negative real number must not silently be interpreted as the real cube root.
The code does not use \code{PowerExpand}.

To establish the precondition for exact root matching, compute the minimal
polynomial of $a$ from the square-free part of the characteristic polynomial
of $M_a$. Indeed,
\[
 \det(TI-M_a)=\operatorname{minpoly}_a(T)^{[L:\Q(a)]}.
\]
Both the known internal $r$ and every candidate in \eqref{eq:phase} are roots
of $\operatorname{minpoly}_a(T^p)$. Thus branch selection is a comparison of
\emph{known roots of a known rational polynomial}, not a test that a residual
is numerically small.

\subsection{Separation-bound comparison}\label{sec:separation}
Let $q$ be a square-free rational polynomial and let $\delta>0$ be a
certified lower bound for separation of its distinct complex roots. Assume
that $u,v$ are already known to be roots of $q$. Evaluate each to absolute
error less than $\epsilon<\delta/8$, producing exact dyadic numbers
$\widetilde u,\widetilde v$. Then
\[
\begin{array}{ll}
 u=v &\Longrightarrow
 |\widetilde u-\widetilde v|<\delta/4,\\
 u\ne v &\Longrightarrow
 |\widetilde u-\widetilde v|>3\delta/4.
\end{array}
\]
Comparison of their squared distance with $\delta^2/4$ therefore decides
equality. Only rational arithmetic is needed after evaluation. The code uses
SymPy's exact Mignotte separation bound and bounded-error evaluation primitive,
with these explicitly separated thresholds; see the underlying polynomial
and root-isolation facilities in \cite{sympypoly}. The precondition is checked by
minimal-polynomial divisibility for arbitrary comparison inputs, or supplied
by exact construction identities for internal calls.

This depends on the correctness of the algebraic and bounded-error backend;
it is not a formally verified interval-arithmetic kernel. Ordinary decimal
approximations are used elsewhere only to order candidate work. They do not
certify equality. The dependency on SymPy 1.14's internal bounded-error API
is explicit and motivates the pinned environment and regression tests.

\section{A separate arithmetic-tower checker}
The output records the prime indices, phase choices, radicands, local
assignments, and final polynomial expression. Besides checking identities
during construction, \code{verify\_tower\_result} independently verifies the
emitted arithmetic tower without constructing any number field or using the
discovered automorphisms.

Let $Z$ represent $\zeta$ and $U_i$ the emitted generators. The checker first
checks the assignment syntax and phase formula, and then reduces the target
polynomial modulo
\begin{equation}\label{eq:ideal}
 \Phi_m(Z),\quad
 U_1^{p_1}-A_1(Z),\quad\ldots,\quad
 U_s^{p_s}-A_s(Z,U_1,\ldots,U_{s-1}).
\end{equation}
Here each $A_i$ is the exact earlier-basis radicand. The root-of-unity phase
disappears upon taking the $p_i$th power. Reduction starts with $U_s$ and
proceeds backwards to $Z$. The relations are monic in their newly introduced
variables, so polynomial remainders suffice. Horner evaluation of $f$ with
reduction after each multiplication avoids unnecessarily large expanded
powers.

If the remainder is zero, the actual emitted radical expression satisfies
$f$ because its assignments satisfy every relation in \eqref{eq:ideal}.
The separation test then identifies it with the selected $\alpha$. This
checker is logically distinct from the splitting-field search. It does not
prove the Galois-group computation correct and it is not a proof object
checked by Lean or another formal kernel; it is an independently executable
exact arithmetic check of a positive answer. On large examples, running the
checker can cost more than constructing the answer.

\section{The complete algorithm and its guarantee}
\begin{lstlisting}
Input: an exact selected algebraic root alpha
f := the irreducible minimal polynomial of alpha over Q
Try finite exact structural shortcuts unless disabled
E := splitting field of f in the selected complex embedding
G := all automorphisms of E/Q, with an exact group table
If the derived series of G is nontrivially stationary:
    return NotSolvableByRadicals
m := product of the distinct prime divisors of |G|
L := E(zeta_m)
H := automorphisms of L fixing zeta_m, by extension tests
Refine H to a series with prime cyclic quotients
For each adjacent pair, unless alpha is already in the fixed field:
    compute a nonzero Kummer eigenvector by rational linear algebra
    optionally improve its radicand with finite trace projections
    express its prime power in the earlier tower basis
    select its exact principal-power phase
    extend the tower basis and append a radical assignment
Express alpha in the final basis and return the radical circuit
\end{lstlisting}

\begin{theorem}[Completeness of the mathematical algorithm]
Assume exact, terminating algorithms for rational and number-field
factorization, primitive elements, field membership, rational linear algebra,
and algebraic root isolation. With no resource bounds, the construction
terminates for every exact algebraic input. It returns a radical expression
for the selected number exactly when such an expression exists.
\end{theorem}
\begin{proof}
The splitting-field stage adjoins finitely many roots of $f$. Automorphism
computation and all subsequent group operations operate on finite sets.
The derived-series criterion correctly distinguishes solvable and
nonsolvable groups. In the solvable case, prime-index refinement terminates
because the group order decreases at every step. The relative-automorphism
construction accounts for the cyclotomic intersection exactly.

Theorem~\ref{thm:kummer} supplies a nonzero generator at every required step.
The basis invariant guarantees each radicand coordinate solve and the final
target coordinate solve. A branch candidate exists and is unique because
all $p_i$th roots of a nonzero number differ by the included roots of unity.
Induction on the assignments identifies each symbolic generator with its
concrete field element; the last coordinate identity therefore gives the
specified $\alpha$. In the nonsolvable case, the Galois criterion excludes
any radical expression.
\end{proof}

The theorem describes an effective algorithm, not an assertion that every
backend call in a particular CAS version is bug-free. The implementation
reports algebraic failures separately and enables budgets by default. It
contains no hard degree-six limit on the general procedure. Conversely,
setting the limits to infinity does not make large instances practical.

\section{Finite shortcuts that matter in practice}
The complete construction is a fallback. A small exact structural identity
can be much more useful than an explicit normal closure.

\subsection{Centered powers and polynomial composition}
After translating away the coefficient of $x^{n-1}$, compute the greatest
common divisor $g$ of the nonzero positive exponents. If $g>1$, the polynomial
is $h(x^g)$. Solve $h$ recursively and take every correlated $g$th-root phase.
For example, $(x+1)^7+2$ immediately gives its real root as
$-1-2^{1/7}$. Exact polynomial decomposition is another route: if
$f=f_1\circ\cdots\circ f_r$, solve successive lower-degree equations,
checking the composition identity before accepting it.

\subsection{Dickson polynomials}
Define
\[
 D_0(x,a)=2,\quad D_1(x,a)=x,\quad
 D_n(x,a)=xD_{n-1}(x,a)-aD_{n-2}(x,a).
\]
Induction proves
\[
 D_n(u+a/u,a)=u^n+(a/u)^n.
\]
Thus $D_n(x,a)=b$ is solved by choosing a nonzero root $v$ of
$v^2-bv+a^n=0$, putting $u=v^{1/n}$, and taking
\[
 x_k=\zeta_n^k u+\frac{a}{\zeta_n^k u},\qquad 0\le k<n.
\]
The two summands are linked. Choosing two $n$th roots independently destroys
the required product and can produce incorrect answers.

\subsection{Reciprocal lifts}
For a monic degree-$2d$ polynomial, test identities of the form
\[
 f(x)=x^d h(x+a/x),\qquad a\in\Q.
\]
The constant term must equal $a^d$. The code enumerates the rational
solutions of this equation, recovers the coefficients of $h$ by triangular
coefficient elimination, and accepts only a zero residual. A root $s$ of $h$
then gives
\[
 x=\frac{s\pm\sqrt{s^2-4a}}2.
\]
Failure to find such a rational $a$ is merely failure of this shortcut.

\subsection{Pair-sum resolvents}
For monic $f$ of degree $n$ with roots $\alpha_1,\ldots,\alpha_n$, put
\[
 S(s)=\prod_{i<j}(s-\alpha_i-\alpha_j).
\]
A direct product expansion gives
\begin{equation}\label{eq:pairs}
 \Res_x(f(x),f(s-x))=2^n f(s/2)S(s)^2.
\end{equation}
The diagonal terms give the first factor; each unordered off-diagonal pair
occurs twice. The code divides out the diagonal factor and checks the square
multiplicities exactly. If a low-degree factor supplies a radical pair sum
$s$, then
\[
 \gcd(f(x),f(s-x))\in\Q(s)[x]
\]
can reveal a small factor containing the desired root. Collisions between
pair sums do not invalidate \eqref{eq:pairs}; multiplicities must be retained.
The implemented search uses input degree at most eight and resolvent factors
of degree at most four. These limits affect this shortcut, not the general
algorithm.

\subsection{Cyclotomic polynomials and real traces}
If $f=\Phi_m$, every root is explicitly $\zeta_m^k$ with $(k,m)=1$. The
code also recognizes real cyclotomic generators by testing whether
$x^{\deg f}f(x+x^{-1})$ is cyclotomic; then roots are
$\zeta_m^k+\zeta_m^{-k}$.

Order recognition is finite. Since
\[
 \frac{\varphi(m)^2}{m}
 =\prod_{p^a\parallel m}p^{a-2}(p-1)^2\ge\frac12,
\]
a cyclotomic polynomial of degree $d$ has $m\le2d^2$. The inequality follows
factor by factor: only $p=2,a=1$ contributes less than one, and that
contribution is $1/2$. An exact cyclotomic test precedes this bounded search.

All these constructions establish that their candidates satisfy the input
polynomial before root matching. Recomputing the minimal polynomial of each
large nested radical is optional, not logically necessary. The low-degree
base cases trust the CAS's exact degree-at-most-four formulas; higher steps
check the stated polynomial identities. The regression report distinguishes
these structural checks from additional minimal-polynomial checks.

\section{The two motivating examples}
\subsection{A sextic that is only a cubic plus a quadratic}
The first example is
\[
 f(x)=x^6+x^4-x^3-x^2-1.
\]
Its structure is visible in the exact identity
\begin{equation}\label{eq:sextic}
 f(x)=x^3\bigl((x-x^{-1})^3+4(x-x^{-1})-1\bigr).
\end{equation}
Set
\[
 c=\left(\frac{9+\sqrt{849}}{18}\right)^{1/3},\qquad
 s=c-\frac{4}{3c}.
\]
The quantity $c$ is positive real. Since $c^3$ solves
$v^2-v-64/27=0$, expansion gives $s^3+4s-1=0$. That cubic has exactly one
real root: its derivative is $3s^2+4>0$. Equation~\eqref{eq:sextic} then
shows that the two real roots of $f$ are
\[
 \frac{s-\sqrt{s^2+4}}2<0,
 \qquad
 \boxed{\alpha=\frac{s+\sqrt{s^2+4}}2>0}.
\]
Hence the boxed expression is the requested second real Wolfram root:
\begin{lstlisting}
Root[-1 - #^2 - #^3 + #^4 + #^6 &, 2]
\end{lstlisting}
No numerical guessing is involved. The other four roots are obtained by
using the other two cubic roots and both quadratic signs, followed by exact
root matching.

The pair-sum route reaches the same intermediate cubic. The computed
resolvent is
\begin{align*}
 S(s)={}&(s^3+4s-1)\bigl(s^{12}-s^9+8s^8-2s^7-s^6\\
       &\hspace{31mm}+4s^5+18s^4+9s^3+2s^2-1\bigr).
\end{align*}
For a root $s$ of the cubic factor, the relevant common factor is
$x^2-sx-1$. This illustrates why the notebook's intermediate-extension ideas
can be effective, while the reciprocal identity is the simpler certificate.

\subsection{The solvable quintic and correlated fifth roots}
The second example is
\[
 P(x)=5x^5-25x^3+25x+6=5D_5(x,1)+6.
\]
Take
\[
 u=\left(\frac{-3+4i}{5}\right)^{1/5}
\]
with its principal value. The five roots are
\[
 x_k=\zeta_5^k u+(\zeta_5^k u)^{-1}.
\]
The largest real root is $x_0=u+u^{-1}$, giving
\begin{lstlisting}
Root[6 + 25 # - 25 #^3 + 5 #^5 &, 5]
  == ((-3 + 4 I)/5)^(1/5) + ((-3 + 4 I)/5)^(-1/5)
\end{lstlisting}
For completeness, write $(-3+4i)/5=e^{i\theta}$ with
$\pi/2<\theta<\pi$. Then $u=e^{i\theta/5}$ and every $x_k$ is real.
Among the five angles modulo $2\pi$, $\theta/5$ is closest to zero, so
$x_0$ is largest. This angle argument explains the ordering; the software
still performs exact selected-root comparison.

\subsection{A genuine degree-five Galois construction}\label{sec:c5}
To exercise the fallback rather than a recognition formula, the tests force
\code{method="galois"} on
\[
 g(x)=x^5+x^4-4x^3-3x^2+3x+1.
\]
Its computed splitting-field group has order five. The cyclotomic
augmentation has rational degree twenty and relative group order five.
With $z=(-1)^{2/5}$, the implemented trace-guided construction for the
largest real root produces
\[
 u=z^3(385z^3+110z^2+220z-66)^{1/5},
\]
and the following polynomial in $u$ and $z$:
\begin{align*}
 \alpha={}&\frac{u^4}{6655}(26z^3+16z^2+41z+6)\\
 &+\frac{u^3}{605}(-10z^3-7z^2+2z-4)\\
 &+\frac{u^2}{55}(4z^3+2z^2+2z+1)+\frac{uz-1}{5}.
\end{align*}
The separate checker verifies $g(\alpha)=0$ by polynomial remainders
modulo $\Phi_5(z)$ and $u^5-(385z^3+110z^2+220z-66)$, then checks the
chosen embedding. The example demonstrates an actual fifth-root Kummer
step, not an invocation of a general quintic radical solver. In automatic
mode the much shorter real-cyclotomic-trace shortcut is preferable.

\section{Using the software}
\subsection{Python interface}
The engine was tested under Python \TestPython{} and SymPy \TestSymPy{}.
Install the pinned dependency and run from the archive root:
\begin{lstlisting}
python -m pip install -r requirements.txt
python tests/test_radical_roots.py
\end{lstlisting}
A library call returns a structured result:
\begin{lstlisting}
import sys
sys.path.insert(0, "python")
import sympy as s
from radical_roots import radicalize, verify_tower_result
x = s.Symbol("x")
r = radicalize(x**6 + x**4 - x**3 - x**2 - 1, index=1)
assert r.status == "Success"
print(r.to_wolfram())
print(r.expanded())
\end{lstlisting}
The Python index is zero-based in SymPy's ordering of the distinct roots.
It must not be obtained by subtracting one from an arbitrary Wolfram complex
root index. The optional rational-rectangle selector avoids this convention
entirely.

The command-line interface accepts a JSON array of rational coefficient
strings, highest degree first, rather than executable Python expressions:
\begin{lstlisting}
python python/radical_roots.py --coefficients "[1,0,1,-1,-1,0,-1]" --index 1 --timeout 300
\end{lstlisting}
Use \code{--method galois} to force the general construction. Defaults are
a field-degree cap of 128 and a 300-second subprocess watchdog. The CLI
uses \code{--max-degree 0 --timeout 0} to disable both; the library uses
\code{max\_degree=None}. Disabling limits can consume substantial time and
memory. The degree cap is checked after constructing an extension, so it is
not a strict bound on transient work or memory.

\subsection{Wolfram Language interface}
Keep the extracted directory structure intact. In Wolfram Language:
\begin{lstlisting}
Get["/full/path/to/RadicalRoots/wolfram/RadicalRoots.wl"];
r = Root[-1 - #^2 - #^3 + #^4 + #^6 &, 2];
a = RootToRadicals[r, "PythonExecutable" -> "/full/path/to/python"];
RootReduce[a - r]
\end{lstlisting}
On Windows, supply the full path to \code{python.exe} when it is not on the
process search path.

The metadata interface is \code{RadicalRepresentation[r]}. To read a saved
Python result, use \code{ReadRadicalJSON[file,r]}; it applies the same Wolfram
verification without launching Python again.

The front end proposes a dyadic rectangle using an approximation, proves
that the selected input lies strictly inside it using exact algebraic
comparisons, and checks its root count with \code{CountRoots}. The documented
complex-rectangle interpretation is used \cite{countroots}. Python independently
rejects boundary roots and requires a unique root in the box. This avoids
assuming compatible complex-root ordering across the two systems.

The returned JSON expression is an arithmetic tree with tags for rationals,
$i$, addition, multiplication, rational powers, and earlier local variables.
It is decoded without evaluating source-code strings. Process invocation
uses an argument list, not a shell command. Finally, \code{RootReduce}
checks exact equality to the original selected value \cite{rootreduce}.
A verification timeout is reported as a resource limit, not as a successful
conversion. The Wolfram integration remains unexecuted in this release;
the supplied \code{Tests.wlt} is intended for validation in a real kernel.

\section{Validation and honest limitations}
The reproducible Python suite reports \TestPassed{} passing tests out of
\TestCount{}. It records the Python and SymPy versions, the source SHA-256,
per-test outcomes, timings, and which answers received an additional
minimal-polynomial check. Galois-Kummer results also receive the separate
triangular-polynomial tower check. Eight additional command-line and
certificate-rejection checks passed and have their own JSON report.

Coverage includes all six conjugates of the motivating sextic, all five
roots of the motivating quintic, real and nonreal pure-cubic roots, quadratic
and biquadratic extensions, a dihedral quartic, and a forced cyclic-quintic
construction. It also includes polynomial composition, a shifted seventh
power, cyclotomic recognition, repeated roots, selected irreducible factors,
a nonsolvable $S_5$ example, explicit resource exhaustion, rejection of
inexact or symbolic coefficients, and rational rectangle selection including
boundary rejection. These are regression cases, not exhaustive certification
of all number fields.

An earlier version that refactored the augmented degree-twenty field hit a
180-second watchdog on the cyclic quintic. The relative-automorphism method
removed that bottleneck, and trace-guided generator selection dramatically
reduced the resulting expression. These observations explain design choices;
they are not a controlled benchmark against Mathematica, SageMath, Magma,
or other computer algebra systems.

The dominant remaining limitation is explicit normal-closure construction.
A degree-$n$ input can have a splitting field of degree up to $n!$, and a
solvable input can still produce large fields and coefficients. Rational
nullspaces, automorphism tables, basis coordinates, and principal-power
verification may all become expensive. Very large notebook examples were
not comprehensively benchmarked, and minimal radical depth or globally
shortest output is not optimized.

This should not be confused with a general complexity impossibility.
Landau and Miller give polynomial-time solvability decision \emph{and}
radical construction with suitable straight-line-program encoding
\cite{landaumiller}. The present explicit-field reference implementation
is not their algorithm and makes no polynomial-time claim. Its advantages
are a direct proof, inspectable exact operations, useful structural shortcuts,
and a clear distinction between a proved negative answer and an unfinished
computation.

\section{Directions for a stronger implementation}
The most valuable next step is to avoid constructing unnecessarily large
normal closures. Imprimitivity blocks, field intersections, and primitive
solvable-group methods can support a Landau--Miller-style architecture.
The arithmetic-tower interface can remain unchanged while discovery is
replaced by a more scalable implementation.

Other improvements are compatible with the existing correctness argument:
reduced integral bases for field arithmetic; lattice reduction and more
systematic Kummer-generator optimization; lazy expression sharing;
parallel independent exact candidate checks; persistent field-factorization
caches; and a formally verified checker for the emitted arithmetic circuit
and root-isolation evidence. An optional search suggested by numerical
relations should only propose candidates, each of which must pass an exact
certificate check before it affects a result.

The conceptual answer to the original question is therefore positive and
algorithmic. The obstacle is not a missing criterion for radical
expressibility. The practical task is to implement its constructive steps,
control field and expression growth, and preserve the selected embedding.
The supplied package addresses each of these explicitly, while keeping its
unexecuted integration layer and finite testing coverage visible.

\begin{thebibliography}{9}
\bibitem{question}
V. Reshetnikov, \emph{Why is ToRadicals not able to handle all cases? Is there
a workaround?}, Mathematica Stack Exchange, question 34011.
\url{https://mathematica.stackexchange.com/questions/34011/}.
The user-supplied \code{FindExtension.nb} was inspected separately.

\bibitem{toradicals}
Wolfram Research, \emph{ToRadicals}, Wolfram Language documentation.
\url{https://reference.wolfram.com/language/ref/ToRadicals.html}.

\bibitem{milne}
J. S. Milne, \emph{Fields and Galois Theory}, version 5.10, 2022.
In particular, Proposition 5.27 and Theorem 5.34.
\url{https://www.jmilne.org/math/CourseNotes/FT.pdf}.

\bibitem{landaumiller}
S. Landau and G. L. Miller, \emph{Solvability by radicals is in polynomial
time}, Journal of Computer and System Sciences \textbf{30} (1985),
no. 2, 179--208.
\url{https://www.cs.cmu.edu/~glmiller/Publications/Papers/LaMi85.pdf}.

\bibitem{sympynf}
SymPy Development Team, \emph{Number Fields}, SymPy 1.14 documentation.
\url{https://docs.sympy.org/latest/modules/polys/numberfields.html}.

\bibitem{sympypoly}
SymPy Development Team, \emph{Polynomial Manipulation Module Reference},
SymPy 1.14 documentation, including exact root comparison and root isolation.
\url{https://docs.sympy.org/latest/modules/polys/reference.html}.
The implementation also uses the 1.14 bounded-error evaluation primitive.

\bibitem{countroots}
Wolfram Research, \emph{CountRoots}, Wolfram Language documentation.
\url{https://reference.wolfram.com/language/ref/CountRoots.html}.

\bibitem{rootreduce}
Wolfram Research, \emph{RootReduce}, Wolfram Language documentation.
\url{https://reference.wolfram.com/language/ref/RootReduce.html}.

\end{thebibliography}
\smallskip
\small Web documentation was consulted on 8 September 2026. The software
pins the actually tested SymPy version; references to current documentation
do not assert compatibility with every future release.
\end{document}
